\chapter{The Business of Market Making}\label{ch:hf:the-business-of-market-making}

In 2025 Virtu Financial, a listed \omterm{def:hf:the-business-of-market-making:emm}{electronic market maker}, reported adjusted net trading income of \$8.6 million a trading day. The same filing shows
what earning it cost: \$770 million of exchange, clearing and order-flow fees, \$528 million of pay for about a thousand employees, \$249 million of
communication and data. Flow Traders, an exchange-traded-fund specialist, earned 2.5 basis points for each euro of the products it traded. This
chapter's first market maker, run for ten simulated hours, earns half a cent a share on the spread and gives almost all of it back within ten seconds. The
business is a small margin on a very large volume, over a cost base that does not shrink when volume does.

\section{What a market maker sells}

A market maker (One Quant Book 1, chapter 1) stands ready to buy at its bid and sell at its ask. What it sells is immediacy: the investor who
crosses the spread trades now instead of waiting for someone who wants the other side. The price of immediacy is half the spread per share, earned
on every fill, and its cost is that some of the people who take it know where the price is going.

\begin{definition}[Electronic market maker]\label{def:hf:the-business-of-market-making:emm}
An \emph{electronic market maker}\index{electronic market maker} is a firm that provides liquidity by algorithm: its quotes, their sizes and their
revisions are set by programs that react to market data in microseconds to milliseconds, across many instruments and venues at once, with humans
supervising the programs rather than the quotes.
\end{definition}

Most \omterm{def:hf:the-business-of-market-making:emm}{electronic market makers} are also high-frequency traders in the legal sense that One Quant Book 10, chapter 24 gives the term (short holding
periods, high message rates, colocated servers); the converse fails, since some high-frequency strategies only take liquidity. Some \omterm{def:hf:the-business-of-market-making:emm}{electronic market
makers} carry obligations, as designated or lead market makers (One Quant Book 1, chapters 19 and 29); most quote because it pays. Four business
lines recur in this book: exchange-traded products and their hedges (chapters 12--14), futures and options on exchanges (chapters 15, 19 and 20), the
dealer markets reached by request for quote or streaming, FX and bonds (chapters 21 and 22), and retail order flow bought from brokers (chapter 23).
Crypto venues (chapters 24 and 25) and event markets (chapter 26) run the same trade on newer infrastructure.

\section{Revenue: capture, volume and volatility}

\begin{definition}[Revenue capture]\label{def:hf:the-business-of-market-making:capture}
A market maker's \emph{revenue capture}\index{revenue capture} is its trading income, net of the fees that scale with trading, per unit traded: per
share, per contract, or per unit of value traded, usually in basis points.
\end{definition}

Revenue is capture times volume. Both move with volatility: more volatile days bring more volume, wider spreads and larger adverse moves, and the
listed firms describe their income as driven by volume and volatility together. Capture itself is what is left of the half-spread after the price
has moved against the fills. One Quant Book 7, chapter 23 splits it exactly.

\begin{proposition}[Where the half-spread goes]\label{prop:hf:the-business-of-market-making:split}
Let fills $i$ have side $s_i=\pm1$, quantity $q_i$, price $p_i$ and time $t_i$; let $m(t)$ be a reference price, $H$ a horizon and $T$ the end of
the period; let $f_i$ be the fee paid on fill $i$ (negative for a rebate). The P\&L of the fills, the final position marked at $m(T)$, is
\begin{align*}
\pnl={}&\underbrace{\sum_i s_iq_i\,(m(t_i)-p_i)}_{\text{spread capture}}
+\underbrace{\sum_i s_iq_i\,\bigl(m(t_i+H)-m(t_i)\bigr)}_{\text{adverse selection}}\\
&+\underbrace{\sum_i s_iq_i\,\bigl(m(T)-m(t_i+H)\bigr)}_{\text{inventory}}-\sum_i f_i .
\end{align*}
\end{proposition}

\begin{proof}
For each fill the three brackets telescope to $m(T)-p_i$; $\sum_i s_iq_i\,(m(T)-p_i)$ is the cash from the fills plus the final position valued at
$m(T)$.
\end{proof}

The identity holds for any reference: the mid, a fair price (chapter 2), or, in a simulation, the efficient price $P^\ast_t$ nobody observes. With
the mid, a passive fill's spread capture is exactly half the spread; the adverse-selection term measures how much of it the next $H$ seconds take
back. Menkveld's study of one large high-frequency market maker on Chi-X and Euronext found the same split in real data: a gross profit of
\euro0.88 per trade, made of \euro1.55 earned on the spread net of fees and a \euro0.68 loss on positions, which earned \euro0.45 when held
under five seconds and lost \euro1.13 when held longer.

\section{Costs: fees, financing, people and machines}

A market maker's costs fall into two groups. Those that scale with trading are exchange and clearing fees (less rebates), payment for order flow
where it buys retail orders, and the financing of its inventory and of the margin behind it. Those that do not scale are people, market data,
connectivity and colocation, servers and exchange memberships. The first group is netted into capture; the second decides how much capture the firm
needs to survive.

\begin{omfigure}
\begin{tikzpicture}
\begin{axis}[width=12cm,height=5.6cm,ybar stacked,bar width=20pt,ymin=0,ymax=4000,ylabel={\small \$ millions},
  xtick={0,1,2,3,4,5,6,7},xticklabels={revenue,fees,interest,staff,data,{opera\-tions},{depre\-ciation},{what remains}},
  xticklabel style={font=\scriptsize,text width=1.5cm,align=center},tick label style={font=\footnotesize},
  table/col sep=comma,enlarge x limits=0.08,grid=major,grid style={gray!20},axis x line*=bottom]
\addplot[draw=none,fill=none] table[x=k,y=bottom]{figdata/hft/01-the-business-of-market-making/waterfall.csv};
\addplot[fill=omThm!50,draw=omThm] table[x=k,y=cost]{figdata/hft/01-the-business-of-market-making/waterfall.csv};
\addplot[fill=omDef!60,draw=omDef] table[x=k,y=total]{figdata/hft/01-the-business-of-market-making/waterfall.csv};
\end{axis}
\end{tikzpicture}

\omcaption{Virtu Financial's 2025 revenue and the costs it paid on the way down, in \$ millions: exchange, clearing and order-flow fees and
interest and dividends expense scale with trading; staff, communication and data, operations and depreciation do not. What remains is before financing interest, one-off items
and tax. Data: Form 10-K for 2025.}\label{fig:hf:the-business-of-market-making:waterfall}
\end{omfigure}

\begin{dated}{2026-09}{Two listed market makers in 2025}\label{dat:hf:the-business-of-market-making:public}
\textbf{Virtu Financial} (Form 10-K for 2025): total revenue \$3,632 million, of which trading income \$2,437 million; brokerage, exchange,
clearance fees and payments for order flow \$770 million; interest and dividends expense \$647 million; employee compensation \$528 million;
communication and data processing \$249 million. Adjusted net trading income (the firm's measure: trading income, commissions and interest and
dividends income, less the two direct cost lines) \$2,145 million, \$8.6 million a day over 248.5 trading days; the market-making segment
\$1,666 million, \$6.7 million a day, 77.7\% of the total. About 1,027 employees in February 2026.

\textbf{Flow Traders} (full-year 2025 results): net trading income \euro485.8 million on \euro1,940 billion of exchange-traded-product value traded;
fixed operating expenses \euro204.1 million (fixed employee \euro97.3 million, technology \euro70.6 million, other \euro36.3 million), variable
employee expenses \euro77.4 million; EBITDA \euro198.9 million, a 41\% margin; 635 full-time employees at the end of the year.
\end{dated}

Three ratios come out of the box. Virtu's fixed-cost lines (staff, communication and data, operations, depreciation) add up to \$940 million, \$3.8
million a trading day against \$8.6 million of net trading income; communication and data alone cost \$1.0 million a trading day. Its adjusted net
trading income was \$2.1 million per employee, against \$0.51 million of pay. Flow Traders' net trading income was 2.5 basis points of its
exchange-traded-product value traded; the ratio overstates what it earned on those products, since the income also includes other trading. Its
fixed costs came to \euro0.32 million per employee, a third of them technology.

\section{What the public record shows about who wins}

The two filings describe firms, not trades. The academic record adds what a trade earns and who earns it. Menkveld's decomposition says the spread
pays and the position costs, so the business is to earn the first and shorten the second. Baron, Brogaard, Hagströmer and Kirilenko ranked
high-frequency traders by relative latency and found that differences in speed account for large differences in performance; firms that moved up the
ranking through a colocation upgrade earned more afterwards, both from short-lived information and from managing risk faster, in market making and
in cross-market arbitrage alike. Speed is therefore not a detail of implementation but part of the product: it decides how often a market maker's
quote is the one a better-informed trader picks off (chapter 9), and how quickly it can hedge.

\begin{remark}[What a filing cannot tell you]\label{rem:hf:the-business-of-market-making:limits}
A filing aggregates thousands of instruments and many strategies. It gives no capture per share for the firm as a whole (Virtu does not report its
volume), no split between passive and aggressive trading, and no view of the losing days inside a year. The rest of this book works from
simulations and from the published research for exactly these reasons, and says so in every caption.
\end{remark}

\section{The shape of the business: fixed costs and scale}

With capture $c$ per unit traded, variable costs $v$ per unit (fees, financing) and fixed costs $F$ per day, a daily volume $V$ earns
\[
\pi(V)=(c-v)\,V-F .
\]

\begin{proposition}[Break-even volume and operating leverage]\label{prop:hf:the-business-of-market-making:leverage}
If $c>v$, profit is zero at $V^\ast=F/(c-v)$, and above it the elasticity of profit to volume is
\[
\frac{d\ln\pi}{d\ln V}=\frac{(c-v)\,V}{(c-v)\,V-F}=\frac{V}{V-V^\ast}>1 .
\]
\end{proposition}

\begin{proof}
Direct differentiation of $\pi(V)$; the second form divides numerator and denominator by $c-v$.
\end{proof}

\begin{example}[A small firm]\label{ex:hf:the-business-of-market-making:small}
A new market maker expects a gross capture of 0.5 basis points of value traded, pays 0.2 basis points in fees and financing, and costs \$40,000 a
trading day to run (about \$10 million a year). It breaks even at \$1.33 billion a day of value traded. At \$2 billion it earns \$20,000 a day,
with an operating leverage of 3: a 10\% fall in volume removes 30\% of its profit. If competition takes capture from 0.5 to 0.4 basis points, the
break-even volume rises to \$2 billion, and the same firm earns nothing (\cref{fig:hf:the-business-of-market-making:breakeven}).
\end{example}

\begin{omfigure}
\begin{tikzpicture}
\begin{axis}[width=11cm,height=5.2cm,xlabel={\small daily value traded (\$ billions)},ylabel={\small profit (\$ thousands a day)},
  tick label style={font=\footnotesize},xmin=0,xmax=3,ymin=-45,ymax=55,grid=major,grid style={gray!25},table/col sep=comma]
\addplot[omThm,thick] table[x=volume_bn,y=profit_k]{figdata/hft/01-the-business-of-market-making/breakeven.csv};
\draw[gray,dashed] (axis cs:0,0) -- (axis cs:3,0);
\draw[omDef,thick] (axis cs:1.333,-45) -- (axis cs:1.333,0);
\node[font=\scriptsize,anchor=west,omDef,fill=white,inner sep=1pt] at (axis cs:1.4,-30) {break-even at \$1.33 billion};
\end{axis}
\end{tikzpicture}

\omcaption{Daily profit of the small firm of \cref{ex:hf:the-business-of-market-making:small} (capture 0.5 basis points, variable costs 0.2,
fixed costs \$40,000 a day) against daily value traded. Data: \texttt{firm.mmecon}.}\label{fig:hf:the-business-of-market-making:breakeven}
\end{omfigure}

The fixed cost base explains three features of the industry that the later chapters meet again. Market makers spread one technology stack over as
many instruments and venues as it can serve, because each new instrument adds volume for little fixed cost. They pay for speed, whose cost is fixed
and whose benefit scales with volume. And the firms that survive a long low-volatility, low-volume period are those with the lowest break-even
volume, not the highest capture.

\section{Tutorial: a first market maker}\label{tut:hf:the-business-of-market-making:tut}

\begin{tutorial}
\textbf{Goal.} Run the simplest market maker in a simulated market and decompose its P\&L. \textbf{End state:} the table below and
\cref{fig:hf:the-business-of-market-making:markout}.
\begin{tutsteps}
\item \textbf{The market.} Book 7's \texttt{firm.tape} simulates one stock at \$100 with a one-cent tick for an hour: liquidity providers, noise
traders whose orders cluster, informed traders who trade towards an efficient price $P^\ast_t$ that jumps at random times, and a news window.
\item \textbf{The harness.} \texttt{firm.mmharness} runs a \emph{Quoter} inside that market: the Quoter's orders are real orders in the book, it
sees the market half a millisecond late and its orders arrive half a millisecond after it sends them. It must quote off the book \emph{without}
its own orders, or it chases them: a version that joined the raw best bid, its own orders included, sent 336,012 messages in the first
1,400 simulated seconds of one run, almost all in its last minute, against 323 for the fixed version. \texttt{Ctx.quote} keeps one order per side at a target price
(\cref{lst:hf:the-business-of-market-making:quote}).
\omcode{code/firm/mmharness/firm_mmharness.py}{135}{147}{Keep one order per side at the target price: cancel the others, send if none is left.}\label{lst:hf:the-business-of-market-making:quote}
\item \textbf{The Quoter.} \texttt{SymmetricQuoter} joins the best bid and ask of the others with 100 shares and stops adding to a position of
500 (\cref{lst:hf:the-business-of-market-making:quoter}). Fees are illustrative: a rebate of 0.1 cent a share on passive fills, a 0.1 cent charge
on aggressive ones.
\omcode{code/firm/mmharness/firm_mmharness.py}{166}{184}{The chapter's first market maker.}\label{lst:hf:the-business-of-market-making:quoter}
\item \textbf{Run and decompose.} \texttt{hf\_business.decomposition(H=10)} runs ten independent hours (seeds 1--10) and splits the P\&L as in
\cref{prop:hf:the-business-of-market-making:split}, against the mid and against the efficient price.
\end{tutsteps}
\end{tutorial}

The Quoter trades 38,270 shares an hour on average, 11\% of the market's volume, and sends 1,256 messages. Its P\&L per hour and per share (ten-hour
means, $H=10$ seconds):

\begin{center}
\small
\begin{tabular}{@{}lrrrrr@{}}
\toprule
reference & spread capture & adverse selection & inventory & fees & total\\
\midrule
mid, \$ per hour & 186.50 & $-177.25$ & $-66.50$ & 38.27 & $-18.98$\\
mid, cents per share & 0.487 & $-0.463$ & $-0.174$ & 0.100 & $-0.050$\\
efficient price, cents per share & $-0.173$ & 0.038 & $-0.013$ & 0.100 & $-0.049$\\
\bottomrule
\end{tabular}
\end{center}

Against the mid, the Quoter earns its half-spread of half a cent on every share and loses almost all of it to adverse selection within ten seconds;
the rebate is what keeps it near zero. Against the efficient price, the spread capture is already negative: on average the Quoter's fills happen
when the efficient price has moved through its quote, so the mid it trades at is stale. The mark-out curve shows the half-spread melting
(\cref{fig:hf:the-business-of-market-making:markout}). The hourly total has a standard deviation of \$89, so the ten-hour mean of $-\$19$ carries a
standard error of \$28: the naive Quoter does not make money, and ten hours cannot say whether it loses.

\begin{omfigure}
\begin{tikzpicture}
\begin{axis}[width=11cm,height=5.2cm,xlabel={\small seconds after the fill},ylabel={\small mark-out (cents a share)},xmode=log,log basis x=10,
  xmin=0.4,xmax=70,ymin=-0.3,ymax=0.55,tick label style={font=\footnotesize},grid=major,grid style={gray!25},table/col sep=comma,
  xtick={0.5,1,2,5,10,20,60},xticklabels={0.5,1,2,5,10,20,60},legend style={font=\scriptsize,at={(0.5,-0.3)},anchor=north},legend columns=2]
\addplot[omDef,thick,mark=*,mark size=1.2pt] table[x=h,y=mid]{figdata/hft/01-the-business-of-market-making/markout.csv};
\addplot[omThm,thick,dashed,mark=square*,mark size=1.2pt] table[x=h,y=truth]{figdata/hft/01-the-business-of-market-making/markout.csv};
\draw[gray] (axis cs:0.4,0) -- (axis cs:70,0);
\legend{against the mid,against the efficient price}
\end{axis}
\end{tikzpicture}

\omcaption{Mark-out of the symmetric Quoter's fills, $s_i\,(m(t_i+h)-p_i)$ in cents a share, against the horizon $h$ (log scale), with $m$ the mid or
the efficient price; quantity-weighted over ten simulated hours. At $h=0$, not drawn on the log scale, the mark-out against the mid is the half-spread, 0.487. Data:
\texttt{hf\_business.markout\_curve}.}\label{fig:hf:the-business-of-market-making:markout}
\end{omfigure}

\textbf{What to change next.} Quote one tick behind the best prices and see volume collapse; switch off the rebate (exercise 7); give the Quoter
a fair price instead of the mid (chapter 2) and a reason to skew (chapter 3).

\section{Build: the harness and the unit economics}\label{bld:hf:the-business-of-market-making:harness}

\begin{build}
\bfield{Purpose} One interface for every market-making strategy of the book, so that a strategy written once runs on Book 7's synthetic market today
and on Book 10's exchange simulator when it lands; and the arithmetic of capture, costs and scale.
\bfield{Interface} \texttt{firm.mmharness}: a \texttt{Quoter} with \texttt{on\_start}, \texttt{on\_market(ctx, t, top)}, \texttt{on\_fill(ctx, fill)}
and optionally \texttt{on\_timer}; \texttt{ctx.send}, \texttt{cancel}, \texttt{take}, \texttt{quote}, \texttt{external}, \texttt{position},
\texttt{cash}, \texttt{working()}; \texttt{run\_tape(quoter, cfg, latency, fees) -> Result} with \texttt{fills}, \texttt{pnl(ref)},
\texttt{decompose(H, ref)} and \texttt{markouts(horizons, ref)}, \texttt{ref} the mid or the efficient price. \texttt{firm.mmecon}:
\texttt{capture\_bp}, \texttt{profit}, \texttt{breakeven\_volume}, \texttt{operating\_leverage}, \texttt{cost\_shares}.
\bfield{Rules} The Quoter never sees the efficient price. Its view of the book excludes its own orders. A cancelled order can still fill while the
cancel travels. Same configuration and seed, same fills.
\bfield{Acceptance tests} \texttt{code/firm/mmharness/tests/}: the decomposition adds up to the marked P\&L against both references; the position
limit holds up to one order in flight; messages stay below a tenth of the market's; an aggressive fill pays the taker fee; runs are deterministic.
\texttt{code/firm/mmecon/tests/}: break-even and operating leverage against their definitions.
\bfield{Stretch} An adapter to \texttt{firm.exchsim} (One Quant Book 10, chapter 26) with a parity test against this one; timers in market time.
\end{build}

\begin{omsources}
\item Virtu Financial, Inc., Form 10-K for the year ended 31 December 2025 (SEC EDGAR).
\item Flow Traders, ``4Q and FY 2025 Results'', 12 February 2026.
\item A.~J.~Menkveld, ``High frequency trading and the new market makers'', \emph{Journal of Financial Markets} 16(4), 2013.
\item M.~Baron, J.~Brogaard, B.~Hagströmer and A.~Kirilenko, ``Risk and return in high-frequency trading'', \emph{Journal of Financial and
Quantitative Analysis} 54(3), 2019.
\end{omsources}

\section{Exercises}

\begin{exercise}[$\star$]\label{exo:hf:the-business-of-market-making:1}
Virtu's adjusted net trading income was \$2,145.3 million over 248.5 trading days. What is it per day?
\end{exercise}

\begin{exercise}[$\star$]\label{exo:hf:the-business-of-market-making:2}
Flow Traders earned \euro485.8 million of net trading income on \euro1,940 billion of exchange-traded-product value traded. Express the ratio in basis
points and say why it overstates the capture on those products.
\end{exercise}

\begin{exercise}[$\star$]\label{exo:hf:the-business-of-market-making:3}
A firm captures 0.5 basis points of value traded, pays 0.2 basis points of variable costs and has \$40,000 of fixed costs a day. What is its
break-even daily value traded?
\end{exercise}

\begin{exercise}[$\star\star$]\label{exo:hf:the-business-of-market-making:4}
The same firm trades \$2 billion a day. Compute its profit and its operating leverage, and its profit after a 10\% fall in volume.
\end{exercise}

\begin{exercise}[$\star\star$]\label{exo:hf:the-business-of-market-making:5}
Against the mid, the symmetric Quoter's spread capture is $+0.487$ cents a share; against the efficient price it is $-0.173$. Explain how both can
be true.
\end{exercise}

\begin{exercise}[$\star\star$]\label{exo:hf:the-business-of-market-making:6}
From \cref{fig:hf:the-business-of-market-making:markout}, between which horizons does the mark-out against the mid cross zero, and what does that
say about how long the Quoter can hold a position?
\end{exercise}

\begin{exercise}[$\star\star\star$]\label{exo:hf:the-business-of-market-making:7}
\emph{Coding.} Rerun the ten hours with no fees at all (\texttt{Fees(0, 0)}). What is the total P\&L per share, and why does nothing else change?
\end{exercise}

\begin{exercise}[$\star\star\star$]\label{exo:hf:the-business-of-market-making:8}
\emph{Find the flaw.} ``We ran our new quoter for an hour in the simulator and it made \$61. It works; let us size it up.''
\end{exercise}

\section{Problem: A Fraction of a Cent}

\begin{problem}[{Weekend problem --- a fraction of a cent}]\label{pb:hf:the-business-of-market-making:1}
A founder wants to start an \omterm{def:hf:the-business-of-market-making:emm}{electronic market maker}. You have the public record, a simulator and a spreadsheet.

\medskip\noindent\textbf{Part I --- The public record.}
\begin{enumerate}
\item Define an \omterm{def:hf:the-business-of-market-making:emm}{electronic market maker} and say what it sells.
\item Give Virtu's 2025 adjusted net trading income per day, and its market-making segment's share of it.
\item Add up Virtu's fixed-cost lines and express them per trading day.
\item Compute Virtu's adjusted net trading income and pay per employee.
\item Compute Flow Traders' net trading income per unit of exchange-traded-product value traded, its fixed costs per employee and its EBITDA margin.
\end{enumerate}

\medskip\noindent\textbf{Part II --- Unit economics.}
\begin{enumerate}[resume]
\item Define \omterm{def:hf:the-business-of-market-making:capture}{revenue capture} and write profit as a function of daily volume.
\item Derive the break-even volume and the operating leverage.
\item For the small firm (0.5, 0.2 basis points, \$40,000 a day), give the break-even volume, and the profit and leverage at \$2 billion a day.
\item What is its profit over 250 trading days at \$2 billion a day?
\item Fixed costs rise to \$50,000 a day. What is the new break-even volume?
\item Capture falls to 0.4 basis points instead. What is the break-even volume now, and the profit at \$2 billion?
\end{enumerate}

\medskip\noindent\textbf{Part III --- The first Quoter.}
\begin{enumerate}[resume]
\item State the P\&L identity of \cref{prop:hf:the-business-of-market-making:split} and prove it.
\item Why must a Quoter quote off the book without its own orders?
\item Give the Quoter's hourly volume, market share and messages.
\item Give its P\&L per share against the mid, part by part.
\item Why is its spread capture negative against the efficient price?
\item How precise is the ten-hour mean of the hourly total?
\end{enumerate}

\medskip\noindent\textbf{Part IV --- The verdict.}
\begin{enumerate}[resume]
\item What did Menkveld and Baron and co-authors find, and what do they imply for this Quoter?
\item State the \emph{named result}: the small firm's break-even daily volume and its operating leverage at \$2 billion a day.
\item In two sentences, what should the founder build first: more capture, more volume, or lower fixed costs?
\end{enumerate}
\end{problem}

\section{Interview questions}

\begin{interviewq}[$\star$ \iqroles{trader}]\label{iq:hf:the-business-of-market-making:1}
What does a market maker sell, and who pays for it?
\end{interviewq}

\begin{interviewq}[$\star$ \iqroles{researcher}]\label{iq:hf:the-business-of-market-making:2}
Decompose a market maker's P\&L into spread capture, adverse selection and inventory. Why does the split depend on a horizon?
\end{interviewq}

\begin{interviewq}[$\star\star$ \iqroles{trader, researcher}]\label{iq:hf:the-business-of-market-making:3}
Your quoter's fills are profitable at the moment of the fill and unprofitable ten seconds later. What is happening, and what do you try first?
\end{interviewq}

\begin{interviewq}[$\star\star$ \iqroles{developer}]\label{iq:hf:the-business-of-market-making:4}
A simulated quoter sends a million messages an hour and fills a few hundred times. What bug do you look for?
\end{interviewq}

\begin{interviewq}[$\star\star$ \iqroles{risk}]\label{iq:hf:the-business-of-market-making:5}
Why is a market maker's profit more volatile than its volume?
\end{interviewq}

\begin{interviewq}[$\star\star\star$ \iqroles{researcher}]\label{iq:hf:the-business-of-market-making:6}
A strategy's hourly P\&L has a mean of $-\$19$ and a standard deviation of \$89. How many independent hours would you need to tell a true mean of
$-\$19$ from zero at two standard errors?
\end{interviewq}
